\documentclass[11pt]{article}
\usepackage{amsmath,amssymb,amsthm}
\newtheorem{theorem}{Theorem}
\newtheorem{corollary}[theorem]{Corollary}
\newcommand{\Qn}{\{0,1\}^n}
\begin{document}
\title{Ordered Root Lists: A Transfer and Its Degree Limitation}
\author{}
\date{November 23, 2026}
\maketitle

This note tests a sufficient inequality mechanism, as opposed to exact
moment cancellation. It does not prove the requested universal head bound.
The obstruction below applies to the stated coupling certificate, not
to arbitrary admissible rational scores.

\section{A reciprocal identity with a positive remainder}

\begin{theorem}[Finite martingale coupling]
Let $u_1,\ldots,u_p,v_1,\ldots,v_q$ be real numbers and let
$\pi_{ij}\ge0$. Write $a_i=\sum_j\pi_{ij}$ and
$b_j=\sum_i\pi_{ij}$, and assume
\[
                  \sum_j\pi_{ij}(v_j-u_i)=0
                  \quad\hbox{for every }i.
\]
If $C+u_i>0$ and $C+v_j>0$ for all indices, then
\[
 \sum_j\frac{b_j}{C+v_j}-\sum_i\frac{a_i}{C+u_i}
 =\sum_{i,j}\frac{\pi_{ij}(v_j-u_i)^2}
                         {(C+u_i)^2(C+v_j)}.
 \tag{1}
\]
Also
\[
 \sum_j b_jv_j^2-\sum_i a_iu_i^2
             =\sum_{i,j}\pi_{ij}(v_j-u_i)^2.
 \tag{2}
\]
Both differences are strictly positive if and only if at least one
$\pi_{ij}>0$ has $u_i\ne v_j$.
\end{theorem}

\begin{proof}
For positive $s,t$, direct multiplication gives
\[
                 \frac1t-\frac1s+\frac{t-s}{s^2}
                         =\frac{(t-s)^2}{s^2t}.
\]
Substitute $s=C+u_i$, $t=C+v_j$, multiply by $\pi_{ij}$,
and sum. The row mean hypothesis cancels the linear term and gives
(1). Expanding the square and using the same hypothesis gives
\[
 \sum_{i,j}\pi_{ij}(v_j-u_i)^2
  =\sum_jb_jv_j^2-2\sum_i a_iu_i^2+\sum_i a_iu_i^2,
\]
which is (2). Every summand on each right side is nonnegative,
and the denominators in (1) are positive. Thus the strictness
conditions are identical.
\end{proof}

\section{The head compiler and the limitation}

Let $u_i(x),v_j(x),a_i(x),b_j(x)$ be affine functions on $\Qn$,
with all $a_i,b_j$ nonnegative there. Suppose that at each vertex
there is a coupling as above, or a coupling in the reverse direction.
The couplings can depend arbitrarily on the vertex and are not required
to be affine. Require strictness at every vertex. Assign label $+1$
to the displayed direction and $-1$ to the reverse direction.

\begin{theorem}[All such labels have threshold degree at most three]
Under these hypotheses, the labels are strictly sign represented by
the polynomial
\[
               Q(x)=\sum_j b_j(x)v_j(x)^2
                            -\sum_i a_i(x)u_i(x)^2,
\]
whose total degree is at most three. If the weights $a_i,b_j$ are
constant, the degree is at most two.

The same labels admit a score with at most $p+q$ admissible heads,
using the affine numerators $b_j$ and $-a_i$ and a common affine
offset in all denominators.
\end{theorem}

\begin{proof}
Apply (2) at each vertex in the stated direction, or in the reverse
direction. Strictness gives $y(x)Q(x)>0$. Each term in $Q$ has
degree at most three, or at most two for constant weights.

To verify the architectural claim explicitly, write all root forms
as $\ell(x)=\ell_0+\sum_k\ell_kx_k$. Choose
\[
 c_k=1+\max_\ell|\ell_k| \quad(0\le k\le n),
             \qquad C(x)=c_0+\sum_{k=1}^n c_kx_k.
\]
Then every $C+\ell$ has a positive constant coefficient and all
variable coefficients strictly positive. It is therefore positive
on the whole cube and admissible in the manuscript model. Formula
(1) shows that
\[
          \sum_j\frac{b_j(x)}{C(x)+v_j(x)}
                -\sum_i\frac{a_i(x)}{C(x)+u_i(x)}
\]
has exactly the prescribed strict signs. The normal-form theorem
implements its $p+q$ fractions as $p+q$ actual heads.
\end{proof}

\begin{corollary}[Majorization does not give a universal construction]
No matter how long the root lists are, the preceding affine-weight
coupling criterion cannot represent parity on $n\ge4$ variables.
For constant weights it cannot represent parity on $n\ge3$ variables.
In particular, replacing moment equalities by these majorization
inequalities does not supply the missing universal construction.
\end{corollary}

\begin{proof}
Put $\chi(x)=(-1)^{x_1+\cdots+x_n}$. For every polynomial $P$ of
degree less than $n$,
\[
                         \sum_{x\in\Qn}\chi(x)P(x)=0.
\]
Indeed every monomial of such a polynomial omits some coordinate;
pairing its values across that coordinate cancels its contribution.
If $\chi(x)P(x)>0$ at every vertex, the displayed sum would be
strictly positive. Thus a polynomial of degree at most three, or
at most two, cannot represent the indicated parity functions.
\end{proof}

For ordinary equal-weight majorization, the coupling is a doubly
stochastic matrix taking the more spread out list to the less spread
out list. The hypothesis is satisfied by that matrix. The preceding
proof uses only its row sums, column sums, and row means, so it does
not rely on a classification theorem for such matrices.

\section{An even stronger ordering requirement}

If two equal-length root lists can instead be matched pointwise with
$u_i(x)\le v_{\sigma(i)}(x)$, with at least one strict inequality,
then
\[
                        L(x)=\sum_jv_j(x)-\sum_i u_i(x)>0.
\]
Reversing the inequalities reverses the sign of $L$. Hence a labeling
certified by one of these two coordinatewise matching directions at
each vertex already has an affine separator. The matching may depend
on the vertex. This stronger ordering condition cannot even certify
two-bit parity.

The analysis does not exclude comparisons valid for the single
reciprocal test function without a convex-order coupling. Establishing
a small-head construction in that unrestricted setting remains the
original unresolved task; no universal feasibility assertion is made
here.
\end{document}
